\chapter{État de l'art}
\label{chap:etat-art}

Avant d'aborder la conception de \csbag, ce chapitre dresse un état de l'art des
technologies, méthodes et bonnes pratiques disponibles à l'entrée en stage. Il
ne décrit pas mes réalisations, présentées au chapitre suivant, mais le cadre
existant dans lequel s'inscrit mon travail~: l'écosystème Microsoft, les
services cloud Azure, l'architecture des applications web modernes, la
conception des bases de données relationnelles et, enfin, les stratégies de
modernisation d'un existant. Chacune de ces notions servira ensuite de grille de
lecture pour analyser et justifier les choix techniques du
chapitre~\ref{chap:travaux}.

\section{L'écosystème Microsoft 365 et SharePoint}
\label{sec:m365}

Le service qui m'a accueilli développe ses solutions autour de Microsoft~365,
une suite d'outils collaboratifs hébergés dans le cloud~\cite{ms_m365}. Au cœur
de cette suite, SharePoint~Online est une plateforme de collaboration et de
gestion documentaire organisée en sites, chaque site regroupant des listes
(données structurées) et des bibliothèques de documents~\cite{ms_sharepoint}.
Dans un groupe fortement décentralisé, ce type de plateforme fournit un socle
commun à des centaines d'entités, ce qui explique sa place centrale au sein du
service.

Plusieurs approches permettent d'accéder par programme aux données de
Microsoft~365 et de SharePoint. Microsoft~Graph constitue l'interface REST
unifiée et recommandée aujourd'hui par Microsoft pour l'ensemble des services
de la suite~\cite{ms_graph}. À côté, deux approches plus anciennes subsistent~:
le modèle objet client (CSOM), spécifique à SharePoint, et l'API REST propre à
SharePoint~\cite{ms_sp_access}.

\begin{table}[H]
  \centering
  \caption{Approches d'accès programmatique à SharePoint~/ Microsoft~365.}
  \label{tab:sp-acces}
  \begin{tabularx}{\textwidth}{>{\bfseries}p{3.4cm} X p{2.9cm}}
    \toprule
    \normalfont\bfseries Approche & \bfseries Description & \bfseries Positionnement \\
    \midrule
    Microsoft Graph & Interface REST unifiée pour tout Microsoft~365
      (utilisateurs, fichiers, sites, etc.) & Approche moderne recommandée \\
    CSOM & Modèle objet côté client, spécifique à SharePoint & Approche
      historique \\
    API REST SharePoint & Accès REST propre à SharePoint & Approche historique \\
    \bottomrule
  \end{tabularx}
\end{table}

\section{Le cloud et les services Azure}
\label{sec:azure}

L'informatique en nuage (\textit{cloud computing}) désigne la fourniture à la
demande de ressources informatiques par un fournisseur, selon un modèle où l'on
consomme un service au lieu d'héberger et d'administrer soi-même
l'infrastructure. La définition de référence, établie par le NIST, distingue
trois modèles de service~: l'infrastructure (IaaS), la plateforme (PaaS) et le
logiciel (SaaS)~\cite{nist2011cloud}. Microsoft~Azure est la plateforme cloud de
Microsoft~; elle propose un ensemble de services managés que l'on assemble selon
les besoins d'une application.

Plusieurs de ces services sont directement pertinents pour une application
d'entreprise. La gestion des identités et de l'authentification est assurée par
Microsoft~Entra~ID, anciennement Azure~Active~Directory, qui prend notamment en
charge l'authentification unique (SSO) via le protocole OpenID~Connect
\cite{ms_entra,ms_oidc}. La protection des données sensibles repose sur
Azure~Key~Vault, un service dédié au stockage centralisé et sécurisé des secrets,
clés et certificats~\cite{ms_keyvault}~; cette approche répond à une bonne
pratique élémentaire, qui consiste à ne jamais inscrire de secret en clair dans
le code source. Enfin, l'exécution de traitements en arrière-plan peut s'appuyer
sur les Azure~WebJobs, déclenchés notamment par une file d'attente
Azure~Storage~Queue, ce qui permet de découpler et de fiabiliser des traitements
asynchrones~\cite{ms_webjobs}.

\begin{figure}[H]
  \centering
  \resizebox{0.95\textwidth}{!}{%
  \begin{tikzpicture}[
    font=\small,
    app/.style={rounded corners=5pt, draw=vinciblue, line width=1pt, fill=vinciblue!6,
      drop shadow={opacity=0.18,shadow xshift=1pt,shadow yshift=-1pt},
      align=center, text width=4cm, minimum height=1.3cm},
    serv/.style={rounded corners=5pt, draw=azurblue!70, line width=0.9pt, fill=white,
      drop shadow={opacity=0.16,shadow xshift=1pt,shadow yshift=-1pt},
      align=center, text width=2.8cm, minimum height=2.4cm, inner sep=4pt},
    conn/.style={-{Stealth[length=2.5mm]}, line width=1pt, azurblue!80}
  ]
    \node[app] (app) {\textbf{Application}\\[2pt]\footnotesize ASP.NET Core + Angular};
    \node[serv, below=2.4cm of app, xshift=-4.9cm] (entra) {};
    \node[serv, right=0.4cm of entra] (kv) {};
    \node[serv, right=0.4cm of kv] (queue) {};
    \node[serv, right=0.4cm of queue] (wj) {};
    % Entra ID : bouclier
    \begin{scope}[shift={([yshift=-0.5cm]entra.north)}]
      \fill[azurblue] (0,0.28) -- (0.22,0.17) -- (0.22,-0.07)
        .. controls (0.22,-0.26) and (0,-0.35) .. (0,-0.35)
        .. controls (0,-0.35) and (-0.22,-0.26) .. (-0.22,-0.07)
        -- (-0.22,0.17) -- cycle;
      \draw[white, line width=1pt] (-0.08,-0.02) -- (-0.01,-0.11) -- (0.1,0.09);
    \end{scope}
    \node[align=center] at ([yshift=-1.5cm]entra.north)
      {\textbf{\footnotesize Entra ID}\\[1pt]\rrole{Authentification \& SSO}\\\rrole{(OpenID Connect)}};
    % Key Vault : clé
    \begin{scope}[shift={([yshift=-0.5cm]kv.north)}]
      \draw[azurblue, line width=1.4pt] (-0.1,0) circle (0.11);
      \draw[azurblue, line width=1.4pt] (0.01,0) -- (0.28,0);
      \draw[azurblue, line width=1.4pt] (0.2,0) -- (0.2,-0.1);
      \draw[azurblue, line width=1.4pt] (0.28,0) -- (0.28,-0.12);
    \end{scope}
    \node[align=center] at ([yshift=-1.5cm]kv.north)
      {\textbf{\footnotesize Key Vault}\\[1pt]\rrole{Secrets, clés}\\\rrole{et certificats}};
    % Queue : messages empilés
    \begin{scope}[shift={([yshift=-0.45cm]queue.north)}]
      \foreach \y in {0.14,0.0,-0.14} {
        \draw[azurblue, line width=0.9pt, rounded corners=1pt] (-0.24,\y-0.045) rectangle (0.24,\y+0.045);}
    \end{scope}
    \node[align=center] at ([yshift=-1.5cm]queue.north)
      {\textbf{\footnotesize Storage Queue}\\[1pt]\rrole{File d'attente}\\\rrole{de messages}};
    % WebJobs : engrenage
    \begin{scope}[shift={([yshift=-0.5cm]wj.north)}]
      \draw[azurblue, line width=1.2pt] (0,0) circle (0.12);
      \foreach \a in {0,45,...,315} \draw[azurblue, line width=1.8pt] (\a:0.12) -- (\a:0.2);
      \fill[azurblue] (0,0) circle (0.04);
    \end{scope}
    \node[align=center] at ([yshift=-1.5cm]wj.north)
      {\textbf{\footnotesize WebJobs}\\[1pt]\rrole{Traitements}\\\rrole{en arrière-plan}};
    % connexions
    \draw[conn] (app.south) .. controls +(0,-0.6) and +(0,0.6) .. (entra.north)
       node[pos=0.5, sloped, above, font=\scriptsize\itshape, text=grisfonce]{s'authentifie};
    \draw[conn] (app.south) .. controls +(0,-0.6) and +(0,0.6) .. (kv.north)
       node[pos=0.55, right, font=\scriptsize\itshape, text=grisfonce]{lit les secrets};
    \draw[conn] (queue.east) -- node[above, font=\scriptsize\itshape, text=grisfonce]{déclenche} (wj.west);
    % cadre
    \begin{scope}[on background layer]
      \node[draw=azurblue!35, line width=0.8pt, rounded corners=6pt, fill=azurblue!4,
            fit=(entra)(wj), inner sep=11pt] (box) {};
    \end{scope}
    \node[anchor=west, font=\footnotesize\bfseries, text=azurblue]
      at ([xshift=5pt,yshift=-2pt]box.north west) {Services Azure managés};
  \end{tikzpicture}%
  }
  \caption{Services Azure managés et leur rôle autour d'une application.}
  \label{fig:azure}
\end{figure}

\section{L'architecture des applications web modernes}
\label{sec:archi-web}

Les applications web ont longtemps reposé sur un modèle multipage, où le serveur
renvoie une page HTML complète à chaque interaction. Le modèle dominant
aujourd'hui est celui de l'application monopage (\textit{Single Page
Application}, SPA)~: une seule page est chargée initialement, puis mise à jour
dynamiquement côté navigateur, le serveur n'échangeant plus que des données. Ces
échanges s'appuient sur une API REST, style d'architecture formalisé par
Fielding, fondé sur la manipulation de ressources via les verbes du protocole
HTTP~\cite{fielding2000rest}. Ce découplage entre un front-end autonome et une
API back-end permet de développer et de faire évoluer les deux indépendamment,
et d'ouvrir l'API à d'autres consommateurs. Dans l'écosystème considéré, ce
modèle se traduit par un front-end Angular~\cite{angular_docs} dialoguant avec
une API ASP.NET~Core~\cite{ms_aspnetcore}.

\begin{figure}[H]
  \centering
  \begin{tikzpicture}[
    font=\footnotesize,
    card/.style={rounded corners=3pt, draw=vinciblue, line width=0.8pt, fill=white,
      align=center, text width=5.2cm, minimum height=0.8cm},
    layer/.style={rounded corners=2pt, draw=vinciblue!55, line width=0.6pt,
      fill=vinciblue!8, align=center, text width=4.8cm, minimum height=0.6cm},
    flow/.style={-{Stealth[length=2.2mm]}, line width=0.9pt, vinciblue},
    tag/.style={rounded corners=2pt, draw=vincired!55, fill=vincired!7,
      font=\scriptsize, text=grisfonce, align=center, inner sep=2pt},
    lbl/.style={font=\scriptsize\itshape, text=grisfonce, fill=white, inner sep=1pt}
  ]
    % Front-end
    \node[card] (spa) {\textbf{Navigateur} — Angular (SPA)};
    % Couches API
    \node[layer, below=1cm of spa] (ctrl) {Contrôleurs (présentation)};
    \node[layer, below=0.2cm of ctrl] (svc) {Services (métier)};
    \node[layer, below=0.2cm of svc] (repo) {Repositories (accès données)};
    \draw[flow] (ctrl) -- (svc);
    \draw[flow] (svc) -- (repo);
    % Cadre API
    \begin{scope}[on background layer]
      \node[draw=vinciblue, line width=0.8pt, rounded corners=4pt, fill=grisclair,
            fit=(ctrl)(svc)(repo), inner sep=8pt] (api) {};
    \end{scope}
    \node[anchor=south west, font=\scriptsize\bfseries, text=vinciblue]
          at ([xshift=2pt,yshift=1pt]api.north west) {API ASP.NET Core};
    % Tags patrons
    \node[tag, right=0.4cm of ctrl] (t1) {DTO};
    \node[tag, right=0.4cm of repo] (t2) {Repository};
    \draw[vincired!55, dotted, line width=0.6pt] (ctrl.east) -- (t1.west);
    \draw[vincired!55, dotted, line width=0.6pt] (repo.east) -- (t2.west);
    % Base de données
    \node[cylinder, shape border rotate=90, aspect=0.3, draw=vincired,
          line width=0.8pt, fill=vincired!8, minimum width=2cm, minimum height=0.9cm,
          align=center, below=1cm of api] (db) {\textbf{SQL Server}};
    % Flèches principales
    \draw[flow] (spa) -- node[lbl,right]{HTTP / REST (JSON)} (api.north);
    \draw[flow] (api.south) -- node[lbl,right]{Entity Framework Core} (db);
  \end{tikzpicture}
  \caption{Architecture en couches d'une application web moderne~: du front-end
    SPA à la base de données relationnelle.}
  \label{fig:archi}
\end{figure}

Côté back-end, les applications d'entreprise s'organisent classiquement en
couches~: une couche de présentation (les contrôleurs qui reçoivent les
requêtes), une couche métier (les services qui portent la logique) et une couche
d'accès aux données. Cette organisation s'accompagne de plusieurs patrons de
conception éprouvés, recensés notamment par Fowler~\cite{fowler2002poeaa}. Le
patron Repository isole l'accès aux données derrière une interface~; l'Unit of
Work regroupe un ensemble de modifications au sein d'une même transaction~; les
objets de transfert de données (DTO) découplent les entités internes de ce que
l'API expose. À cela s'ajoute l'injection de dépendances, principe qui consiste
à fournir à un composant les services dont il a besoin plutôt que de les laisser
les instancier lui-même, ce qui améliore la testabilité et le découplage
\cite{seemann2019di,ms_di}.

L'accès aux données relationnelles s'appuie fréquemment sur un
\textit{mapping} objet-relationnel (ORM), qui traduit automatiquement entre les
objets du langage et les tables de la base. Entity~Framework~Core est l'ORM de
référence de l'écosystème .NET~\cite{ms_efcore}. Un ORM accélère le
développement et réduit le code répétitif, au prix d'un contrôle plus fin qu'offre
en contrepartie l'accès direct par requêtes (ADO.NET), parfois préféré pour les
traitements les plus sensibles à la performance.

\section{Les bases de données relationnelles et la normalisation}
\label{sec:normalisation}

Le modèle relationnel, introduit par Codd, organise les données en relations
(tables) reliées entre elles, et constitue le fondement des systèmes de gestion
de bases de données actuels~\cite{codd1970relational}. La qualité d'un schéma
relationnel repose en grande partie sur sa normalisation, c'est-à-dire son
organisation en formes normales successives (1NF, 2NF, 3NF) visant à éliminer la
redondance des données~\cite{codd1972normalization}.

L'enjeu de la normalisation est concret. Lorsqu'une même information est dupliquée
en plusieurs endroits, par exemple une valeur saisie en texte libre répétée sur
de nombreuses lignes, la base s'expose à des anomalies de mise à jour~: une
correction partielle laisse des incohérences, et les fautes de saisie créent de
faux doublons~\cite{elmasri2016fundamentals}. La normalisation répond à ce
problème en extrayant les données répétées dans des tables de référence
dédiées. Les relations de type plusieurs-à-plusieurs sont alors modélisées par
des tables d'association, qui relient deux entités et peuvent porter des
attributs propres à la relation~\cite{silberschatz2019database}. Ces principes
constituent la grille de lecture qui servira, au chapitre~\ref{chap:travaux}, à
analyser le schéma existant et à justifier son évolution.

\section{Dette technique et modernisation applicative}
\label{sec:dette-technique}

Toute application vit et évolue. Les lois de l'évolution logicielle énoncées par
Lehman rappellent qu'un logiciel utilisé doit être continuellement adapté, sous
peine de voir son utilité décroître, et que sa complexité tend à augmenter avec
le temps si aucun effort n'est consenti pour la maîtriser~\cite{lehman1980evolution}.
Cette dégradation progressive est captée par la notion de dette technique,
métaphore introduite par Cunningham~: privilégier une solution rapide à une
solution propre revient à contracter une dette, qu'il faudra rembourser plus tard
sous forme d'efforts de maintenance accrus~\cite{cunningham1992wycash}.

Rembourser cette dette passe par le \textit{refactoring}, discipline formalisée
par Fowler, qui consiste à améliorer la structure interne d'un code sans en
modifier le comportement observable~\cite{fowler2018refactoring}. À l'échelle
d'une application entière, deux stratégies de modernisation s'opposent~: la
réécriture complète, qui repart de zéro mais concentre les risques, et la refonte
incrémentale, qui fait évoluer l'existant par étapes tout en le gardant
fonctionnel. Le patron \textit{Strangler~Fig} illustre cette seconde voie, où le
nouveau système remplace progressivement l'ancien~\cite{fowler2004strangler}.
Cette démarche progressive s'inscrit dans le cadre plus large de la maintenance
logicielle, notamment évolutive, telle que la décrivent les référentiels du génie
logiciel~\cite{iso14764}.

Ces bonnes pratiques, issues de la littérature, fournissent le cadre théorique
sur lequel s'appuie la suite du rapport. Le chapitre suivant confronte cet état
de l'art à l'existant de \csbag et présente, à sa lumière, la conception et la
réalisation de sa version~2.